\documentclass{article}

% Language setting
% Replace `english' with e.g. `spanish' to change the document language
\usepackage[english]{babel}

% Set page size and margins
% Replace `letterpaper' with`a4paper' for UK/EU standard size
\usepackage[letterpaper,top=2cm,bottom=2cm,left=3cm,right=3cm,marginparwidth=1.75cm]{geometry}

% Useful packages
\usepackage{amsmath}
\usepackage{graphicx}
\usepackage[colorlinks=true, allcolors=blue]{hyperref}
\usepackage{indentfirst}
\usepackage{enumitem}
\usepackage{physics}

\title{Exercício 4  -  Cálculo Numérico para Geociência}


\begin{document}
\maketitle

Nesta atividade buscaremos retornar aos conceitos apresentados durante o curso, para que possamos revisar os conceitos de mínimos quadrados.

A temperatura da crosta terrestre é dada pela seguinte expressão
\begin{align}
    \label{eq: EDO_crosta}
    \frac{d^2T(z)}{d^2z} = -\frac{H}{k},
\end{align}
sendo $T$ a temperatura medida em uma profundidade $z$, $H$ é a geração de calor radiogênico e $k$ é a condutividade da rocha. Para esse problema, assumiremos $H$ e $k$ como constantes em função da profundidade para  deixar o problema mais simples.

Podemos resolver de forma que obtemos a seguinte expressão para a temperatura em função da profundidade
\begin{align}
    T(z) = T_0 + \qty(\frac{q_0}{k})z + \qty(\frac{-H}{2k})z^2,
\end{align}
sendo $q_0$ o fluxo de calor com a superfície e $T_0$ a temperatura da superfície.


Com a intenção de verificar esse modelo, cientistas tomaram o seguinte conjunto de dados
\begin{table}[h!]
    \centering
    \begin{tabular}{|c|c|c|}
       \hline
       z(km)& T(z)(°C) & $\sigma_T$\\
       \hline
       -1 & 35 & 2\\
       -2 & 56 & 1\\
       -5 & 95 & 4\\
       -8 & 180& 2\\
       -10& 210& 3\\
       -12& 248& 2\\
       \hline
    \end{tabular}
    \label{tab:placeholder}
\end{table}

E com esses dados precisamos agora determinar quais os valores dos coeficientes de nosso modelo.


1) Use MMQ para encontrar $T_0 $,$ \qty(\frac{q_0}{k}) $ e $\qty(\frac{H}{2k})$. Considere a incerteza para definir os pesos dos dados, de acordo com o que foi ensinado na aula de 11.06. Use 3 algarismos significativos sempre que anotar os dados. Note: a coordenada z está orientada de baixo para cima, e estamos descendo em z. Então TOME CUIDADO COM SINAIS!!!

2) O resultado $T_0$ faz sentido físico?

3) Feito o ajuste da curva, determine o valor do coeficiente de determinação, $R^2$, desse ajuste e justifique sua qualidade.

\end{document}